\setcounter{page}{35}
\section{\S4. Quasi-isomorphisms and the Derived Category}

Let \(\cat{A}\) be an abelian category, and let \(K(\cat{A})\) be the triangulated category of complexes up to homotopy introduced in \S2. A morphism \(f\colon X^\bullet\to Y^\bullet\) in \(K(\cat{A})\) is called a \textit{quasi-isomorphism} if it induces an isomorphism on all cohomology groups. Denote by \(\mathsf{Qis}\) the collection of all quasi-isomorphisms in \(K(\cat{A})\).

\medskip
\textbf{Proposition 4.1.}
The set \(\mathsf{Qis}\) is a multiplicative system in \(K(\cat{A})\).

\textit{Proof.} This follows from the more general result below.

\medskip
\textbf{Proposition 4.2.}
Let \(\cat{C}\) be a triangulated category, \(\cat{A}\) an abelian category, and \(H\colon\cat{C}\to\cat{A}\) a cohomological functor. Let \(S\) be the class of morphisms \(s\in\cat{C}\) such that \(H(T^i(s))\) is an isomorphism for all \(i\in\mathbb{Z}\). Then \(S\) is a multiplicative system in \(\cat{C}\) and is compatible with the triangulation.

\textit{Proof.} We verify axioms \(\text{FR1}\)--\(\text{FR5}\). Axioms \(\text{FR1}\) and \(\text{FR4}\) are trivial. Axiom \(\text{FR5}\) follows from the long exact sequence of a cohomological functor and the five-lemma.

\setcounter{page}{36}
To prove \(\text{FR2}\), consider a diagram
\[
X\xrightarrow{u} Y \xleftarrow{s} Z,\qquad s\in S.
\]
Complete \(s\) to a triangle \((Z,Y,N,s,f,g)\). Complete \(fu\) to a triangle \((W,X,N,t,fu,h)\). Since \((u,\id_N)\) maps two sides of the second triangle to the first, there exists a morphism \(v\colon W\to Z\) giving a morphism of triangles, satisfying \(sv = ut\).

It remains to show \(t\in S\). Since \(s\in S\), the long exact sequence of the first triangle implies \(H(T^i(N))=0\) for all \(i\in\mathbb{Z}\). Substituting into the long exact sequence of the second triangle yields that \(H(T^i(t))\) is an isomorphism for all \(i\in\mathbb{Z}\), so \(t\in S\). The dual statement of \(\text{FR2}\) is proved similarly.

\setcounter{page}{37}
For \(\text{FR3}\), consider \(f,g\colon X\to Y\). It suffices to prove equivalence of:
\begin{enumerate}
\item[(i)] \(\exists\;s\colon Y\to Y'\in S\) with \(sf=sg\);
\item[(ii)] \(\exists\;t\colon X'\to X\in S\) with \(ft=gt\).
\end{enumerate}
It is enough to treat the case \(g=0\), so we need:
(i) \(\exists\,s\colon Y\to Y'\in S,\ sf=0\)
\(\iff\)
(ii) \(\exists\,t\colon X'\to X\in S,\ ft=0\).

Assume (i). By \(\text{TR1},\text{TR2}\), form a triangle \((Z,Y,Y',v,s,u)\). Since \(sf=0\), Proposition 1.1(b) gives \(g\colon X\to Z\) with \(f=vg\). Form a triangle \((X',X,Z,t,g,w)\). Again by Proposition 1.1(b), \(ft=0\). Since \(s\in S\), \(H(T^i(Z))=0\) for all \(i\), which implies \(t\in S\). The converse (ii)\(\Rightarrow\)(i) is symmetric.

\medskip
\textbf{Definition.}
The \textit{derived category} of \(\cat{A}\) is
\[
D(\cat{A}) := K(\cat{A})_{\mathsf{Qis}}.
\]
Similarly define
\[
D^+(\cat{A})=K^+(\cat{A})_{\mathsf{Qis}},\quad
D^-(\cat{A})=K^-(\cat{A})_{\mathsf{Qis}},\quad
D^b(\cat{A})=K^b(\cat{A})_{\mathsf{Qis}}.
\]
By Proposition 3.3, these are full subcategories of \(D(\cat{A})\), and
\[
D^+(\cat{A})\cap D^-(\cat{A}) = D^b(\cat{A}).
\]

\setcounter{page}{38}
\textbf{Definition.}
Let \(\cat{A}\) be abelian, and \(\cat{A}'\subset\cat{A}\) a \textit{thick} full abelian subcategory (closed under extensions in \(\cat{A}\)). Define \(K_{\cat{A}'}(\cat{A})\) to be the full subcategory of \(K(\cat{A})\) consisting of complexes \(X^\bullet\) whose cohomology objects \(H^i(X^\bullet)\) lie in \(\cat{A}'\). Since \(\cat{A}'\) is thick, \(K_{\cat{A}'}(\cat{A})\) is a triangulated subcategory of \(K(\cat{A})\).

Define
\[
D_{\cat{A}'}(\cat{A}) := K_{\cat{A}'}(\cat{A})_{\mathsf{Qis}}.
\]
By Proposition 3.3, \(D_{\cat{A}'}(\cat{A})\) is the full subcategory of \(D(\cat{A})\) of complexes with all cohomology in \(\cat{A}'\). One defines similarly \(K_{\cat{A}'}^+(\cat{A}), D_{\cat{A}'}^+(\cat{A})\), etc., for bounded-below complexes.

\textbf{Remark.} There is a natural functor \(D(\cat{A}')\to D_{\cat{A}'}(\cat{A})\), which is in general neither injective nor surjective (cf.\ Proposition 4.8).

\medskip
\textbf{Example.}
A morphism \(f\colon X^\bullet\to Y^\bullet\) of complexes induces the zero morphism in \(D(\cat{A})\) if and only if:

\((*)\) There exists \(s\colon Y^\bullet\to Y'^\bullet\in\mathsf{Qis}\) such that \(sf\) is homotopic to zero (equivalently, \(\exists\;t\colon X'^\bullet\to X^\bullet\in\mathsf{Qis}\) such that \(ft\) is homotopic to zero).

\setcounter{page}{39}
\begin{enumerate}
\item If \(f\) is homotopic to zero, then \((*)\) holds. The converse is false.
Take the complex
\[
X^\bullet:\; 0\to\mathbb{Z}\xrightarrow{2}\mathbb{Z}\to\mathbb{Z}/2\mathbb{Z}\to 0,
\]
and \(f=\id_{X^\bullet}\). Let \(g\colon X^\bullet\to 0\) be the zero map; then \(g\in\mathsf{Qis}\) and \(gf=0\), but \(\id_{X^\bullet}\) is not homotopic to zero.

\item If \(f\) is homotopic to zero, it induces the zero map on cohomology; the converse is false.
Take
\[
X^\bullet,Y^\bullet:\quad 0\to\mathbb{Z}\xrightarrow{2}\mathbb{Z}\to 0.
\]
The identity map induces zero on cohomology, but there is no \(t\in\mathsf{Qis}\) with \(ft\simeq 0\). Any homotopy operator would force \(2k(x)=1\), impossible in \(\mathbb{Z}\).
\end{enumerate}

\setcounter{page}{40}
For morphisms \(f,g\colon X^\bullet\to Y^\bullet\) of complexes, the implications below are all strict:
\[
f=g \;\Longrightarrow\; f\simeq g \;\Longrightarrow\; f=g \text{ in }D(\cat{A}) \;\Longrightarrow\; H^i(f)=H^i(g)\;\forall i.
\]

\medskip
\textbf{Proposition 4.3.}
Let \(\cat{A}\) be identified with the full subcategory of \(D(\cat{A})\) consisting of complexes concentrated in degree zero. Then the natural embedding \(\cat{A}\hookrightarrow D(\cat{A})\) is an equivalence of categories.

\textit{Proof.} Left to the reader.

We now prove three lemmas, which give an alternative description of \(D^+(\cat{A})\) when \(\cat{A}\) has enough injectives.

\medskip
\textbf{Lemma 4.4.}
Let \(\cat{A}\) be abelian, \(f\colon Z^\bullet\to I^\bullet\) a morphism of complexes. Assume:
\begin{enumerate}
\item[(1)] \(Z^\bullet\) is acyclic;
\item[(2)] Each \(I^p\) is injective in \(\cat{A}\);
\item[(3)] \(I^\bullet\) is bounded below.
\end{enumerate}
Then \(f\) is homotopic to zero.

\textit{Proof.} Standard and straightforward.

\setcounter{page}{41}
\medskip
\textbf{Lemma 4.5.}
Let \(\cat{A}\) be abelian, \(s\colon I^\bullet\to Y^\bullet\) a morphism of complexes. Assume:
\begin{enumerate}
\item[(1)] \(s\) is a quasi-isomorphism;
\item[(2)] Each \(I^p\) is injective;
\item[(3)] \(I^\bullet\) is bounded below.
\end{enumerate}
Then \(s\) admits a homotopy inverse.

\textit{Proof.} Form the mapping cone \(Z^\bullet = T(I^\bullet)\oplus Y^\bullet\) of \(s\). Since \(s\) is a quasi-isomorphism, \(Z^\bullet\) is acyclic. The natural map \(v\colon Z^\bullet\to T(I^\bullet)\) satisfies the hypotheses of Lemma 4.4, hence is homotopic to zero. Let \((k,t)\colon T(I^\bullet)\oplus Y^\bullet\to I^\bullet\) be the corresponding homotopy operator. Then
\[
v = (\id_{I^\bullet},0) = (k,t)\circ d_Z + d_I\circ(k,t).
\]
Separating components yields the required homotopy inverse.